% !TITLE 山居秋暝
% !AUTHOR 王维
% !DYNASTY 唐
% !GENRE 五言律诗
% !ORDER 40

\begin{poem}
空山新雨后，天气晚来秋。
明月松间照，清泉石上流。
竹喧归浣女\textsuperscript{1}，莲动下渔舟。
随意春芳歇\textsuperscript{2}，王孙自可留\textsuperscript{3}。
\end{poem}

\begin{annotations}
\notetext{1}{竹喧归浣女：竹林中传来喧笑声，是洗衣女子归来了。浣（huàn），洗濯。王维写人不写面目，只写声音和动作。}
\notetext{2}{春芳歇：春天的芳华凋谢了。随意，随它去吧。春天过去了没关系，秋天也有秋天的美。}
\notetext{3}{王孙自可留：王孙（贵族子弟，泛指君子）自然可以留在这里。王维自比——这山中值得留下，他不是无处可去，是主动选择留在这里。}
\end{annotations}

\begin{appreciation}
这首诗写于王维隐居辋川时期，是他山水诗的代表作之一。暝（míng）是日落、傍晚的意思。秋暝，就是秋天的黄昏。

空山新雨后，天气晚来秋——空旷的山林刚下过一场雨，傍晚时分，秋意渐浓。空山，不是空无一人的山，是远离尘嚣的山。新雨后，空气清新，万物如洗，整个世界焕然一新。晚来秋，秋天在傍晚时分显得格外深沉。

明月松间照，清泉石上流——明亮的月光透过松枝洒下，清澈的泉水在石头上流淌。这是全诗最经典的两句，也是王维诗中有画的最佳例证。月光是静的，泉水是动的；松间是幽暗的，石上是有光的。一明一暗，一动一静，构成了一个完美的画面。照字写月光如流水般倾泻，流字写泉水如月光般清澈。两个字互文见义，让人分不清哪里是月光，哪里是泉水。

竹喧归浣女，莲动下渔舟——竹林中传来喧闹声，是洗衣的女子回来了；莲花摇动，是渔船顺流而下。浣（huàn）是洗涤。这两句从静转到了动，从景转到了人。但王维写人，不写面目，只写声音和动作：竹喧、莲动。你听见竹林里的笑声，看见莲花在摇晃，便知道有人来了。这种写法，比直接写人更有韵味。

随意春芳歇，王孙自可留——春天的芳华随意消歇吧，王孙自然可以留在这里。王孙原指贵族子弟，这里王维自比。春芳歇，是春天的花谢了，但秋天有秋天的美。自可留，是诗人主动选择留下，不是因为无处可去，是因为这里值得留下。

这首诗像一幅长卷：从空山新雨，到明月清泉，再到竹喧莲动，最后定格在王孙自留。每一帧都是画，连起来就是王维的隐居生活。他不是在逃避世界，是在寻找一种更本真的生活方式。在那个世界里，有月光，有泉水，有浣女的笑声，有渔舟的轻摇——有这些，就够了。
\end{appreciation}
